\subsection{Exact budgets and permutation symmetry}
\label{sec:tie-symmetry}
The particle-ID tie rule is a limitation of the reference controller. To state
the obstruction precisely, let $X$ contain all information permitted to the
selector, including candidate pairs, scores, and any history or coordinates
used for tie breaking. Relabeling particles acts on $X$ and on unordered pairs.
Assume candidate construction is equivariant and the budget $B$ is invariant.
Let $G_X=\{\pi:\pi X=X\}$ be the stabilizer of the full input, with orbits
$O_1,\ldots,O_m$ on the optional pairs.

\begin{proposition}
A deterministic exact-$B$ selector that is equivariant on the relabeling orbit
of $X$ exists if and only if $B$ is a sum of sizes of a subset of the $O_j$.
\end{proposition}
\begin{proof}
Equivariance requires $S(X)=S(gX)=gS(X)$ for every $g\in G_X$, so $S(X)$ must
be a union of stabilizer orbits. Conversely, choose such a union $S_0$ at a
representative $X_0$ and define $S(\pi X_0)=\pi S_0$. If
$\pi X_0=\rho X_0$, then $\rho^{-1}\pi\in G_{X_0}$ fixes $S_0$, making the
definition independent of the chosen relabeling.
\end{proof}

For example, four optional spokes incident to a center, with indistinguishable
leaves under every permitted input feature, form one orbit of size four. An
exact $25\%$ budget requires one spoke and cannot be deterministic and
equivariant on that input. Distinct coordinate features can destroy this
stabilizer; a visually symmetric drawing alone is insufficient.

Equal pair priorities do not by themselves imply this impossibility.
The priority $\max(s_i,s_j)$ can give identical values to several pairs incident
to a node whose score exceeds those neighbors' scores, even when all node scores are distinct. On that domain,
ordering by $(\max(s_i,s_j),\min(s_i,s_j))$ uniquely orders unordered pairs
equivariantly while preserving the primary objective. Thus the frozen ID tie
rule can break equivariance even when a deterministic alternative exists.

Alternatively, iid continuous keys, independent of the selector input and
attached to unordered pairs, can break only cutoff ties. This gives an exact budget and equivariance in
distribution. If keys are transported with the relabeling, corresponding
pairs have identical primary and secondary priorities, yielding pathwise
equivariance almost surely. Restarting a scalar PRNG seed on a reordered pair
array does not implement that coupling. These alternatives were not installed
in the frozen experiment; no empirical effect of the tie rule on accuracy or
runtime has been established here.
